\begin{recipe}{Bruschetta}
\kicker[Voorgerecht,Vegetarisch,Italiaans]{Voorgerecht · vegetarisch · Italiaans}
\index[register]{Bruschetta}
\index[register]{Tomaat}
\index[register]{Knoflook}
\index[register]{Basilicum}

\meta{20 minuten}

\lettrine{B}{ruschetta} eten we wel eens op de camping. Het is een authentiek
Italiaans gerecht en meteen een goede manier om oudbakken brood een goede
bestemming te geven.

\blockrule{Ingrediënten}
\begin{ingredients}
  \ing{een paar sneden}{oudbakken brood (bijvoorbeeld zelfgemaakte ciabatta,
        zie p.~\pageref{recipe:ciabatta})}
  \ing{4}{rijpe, stevige roma tomaten}\index[register]{Tomaat}
  \ing{1 teen}{knoflook}\index[register]{Knoflook}
  \ing{1 flinke hand}{verse basilicum, grof gescheurd}\index[register]{Basilicum}
  \ingb{olijfolie}
  \ingb{zeezout en versgemalen zwarte peper, naar smaak}
\end{ingredients}

\blockrule{Bereiding}
\begin{steps}
  \item Was en snijd de tomaten in vieren. Verwijder de natte zaadlijsten om
        zompig brood te voorkomen en snijd het vruchtvlees in kleine blokjes.
  \item Meng de tomatenblokjes in een kom met de basilicum, een flinke scheut
        olijfolie, peper en zout. Laat dit 10 minuten intrekken op
        kamertemperatuur.
  \item Rooster de sneden brood, ingesmeerd met een beetje olijfolie, op
        middelhoog vuur in een koekenpan goudbruin en knapperig, en draai ze
        op tijd om. Dit kan ook onder de grill in de oven.
  \item Wrijf het warme, krokante brood direct in met het rauwe teentje
        knoflook: door de ruwe structuur van het brood raspt de knoflook
        vanzelf.
  \item Besprenkel het brood met een beetje olijfolie, schep het
        tomatenmengsel erop en serveer direct.
\end{steps}
\end{recipe}
